%% Question bank chapter file - generated by bank_tools.py build (do not edit while the bank is still being built:
%% it is regenerated from the per-test data in Question Bank\data). Every question is one bankq block:
%%   \begin{bankq}[id=..., source={...}, paper=.., qno=.., type=..., subtopic={...}, answer={...}, answer_src=key|solved,
%%                 working={...}, difficulty=1-10, diff_parts={K S I C F}, diff_note={...},
%%                 group=<passage id or empty>, also_in={duplicate copies in other tests}]
%% difficulty = min(10, K+S+I+C+F-1): K concepts 1-3, S steps 1-4, I insight 0-2, C computation 0-2, F format 0-1
%%   ...question LaTeX, options, figures (\bfig / \bfigi -> img/)...
%%   \end{bankq}
%% Passages shared by several questions are bankpassage blocks; their questions carry group=<passage id>.
\documentclass[10pt]{article}
\usepackage[a4paper,margin=1.8cm]{geometry}
\usepackage[T1]{fontenc}
\usepackage{amsmath,amssymb,mathtools}
\usepackage{graphicx}
\usepackage[export]{adjustbox}
\usepackage{xcolor,needspace,keyval}
\usepackage{enumitem}
\usepackage{booktabs,array,tabularx}
\usepackage[hidelinks]{hyperref}
\setlength{\parindent}{0pt}
\setlength{\parskip}{3pt}
\makeatletter
\define@key{bq}{id}{\def\bq@id{#1}}\define@key{bq}{source}{\def\bq@src{#1}}\define@key{bq}{paper}{\def\bq@paper{#1}}
\define@key{bq}{qno}{\def\bq@qno{#1}}\define@key{bq}{type}{\def\bq@type{#1}}\define@key{bq}{subtopic}{\def\bq@sub{#1}}
\define@key{bq}{answer}{\def\bq@ans{#1}}\define@key{bq}{answer_src}{\def\bq@asrc{#1}}\define@key{bq}{working}{\def\bq@work{#1}}
\define@key{bq}{group}{\def\bq@group{#1}}\define@key{bq}{also_in}{\def\bq@also{#1}}
\define@key{bq}{difficulty}{\def\bq@diff{#1}}\define@key{bq}{diff_parts}{\def\bq@dparts{#1}}\define@key{bq}{diff_note}{\def\bq@dnote{#1}}
\newenvironment{bankq}[1][]{\def\bq@work{}\def\bq@also{}\def\bq@group{}\def\bq@diff{?}\def\bq@dparts{}\def\bq@dnote{}\setkeys{bq}{#1}%
  \par\needspace{4\baselineskip}\medskip\noindent{\small\textbf{\bq@id}\hfill\color{gray}[\bq@src, Paper \bq@paper, Q\bq@qno]
  \bq@type\ $\cdot$ \bq@sub\ $\cdot$ \textbf{D\bq@diff}}\par\nopagebreak\relax}
  {\par{\small\color{teal}Answer: \textbf{\bq@ans}\ifx\bq@work\@empty\else\ (solved: \bq@work)\fi
  \ifx\bq@also\@empty\else\par\color{gray}Also in: \bq@also\fi}\par}
\define@key{bp}{id}{\def\bp@id{#1}}\define@key{bp}{source}{\def\bp@src{#1}}\define@key{bp}{questions}{\def\bp@qs{#1}}
\newenvironment{bankpassage}[1][]{\setkeys{bp}{#1}\par\needspace{4\baselineskip}\bigskip\noindent
  {\small\textbf{Passage \bp@id}\hfill\color{gray}[\bp@src; questions \bp@qs]}\par\nopagebreak\relax\itshape}{\par}
\makeatother
\newcommand{\bfig}[2]{\par\smallskip{\centering\includegraphics[width=#1pt,max width=0.62\linewidth,max height=0.32\textheight]{img/#2}\par}\smallskip}
\newcommand{\bfigi}[2]{\includegraphics[width=#1pt,max width=\linewidth,max height=0.32\textheight,valign=c]{img/#2}}
\newcommand{\bmissing}[1]{\textit{[figure missing in the source conversion: \detokenize{#1}]}}
\newenvironment{qoptions}{\begin{itemize}[leftmargin=1.8em,itemsep=1pt,topsep=2pt,label={}]}{\end{itemize}}
\newcommand{\opt}[2]{\item[(#1)] #2}
\providecommand{\fwstart}{}\providecommand{\fwend}{}\providecommand{\qfullsplit}{}
\newcolumntype{L}[1]{>{\raggedright\arraybackslash}p{#1}}
\newcolumntype{C}[1]{>{\centering\arraybackslash}p{#1}}
\newcolumntype{R}[1]{>{\raggedleft\arraybackslash}p{#1}}
\providecommand{\sectionhdr}[1]{}\providecommand{\subsectionhdr}[1]{}


\begin{document}
{\Large\bfseries Definite Integration --- Question Bank}\par\smallskip
{\small Mathematics chapter M27. Questions from all converted test papers, grouped by source series. Answers: printed key where it matches the paper, otherwise solved (answer\_src=solved, with working). Exact repeats across papers appear once, with also\_in listing the other copies.}\par\medskip\hrule

\section*{Score-I papers}
\begin{bankq}[id=SI12-P1-MATH-S1-Q04, source={Score-I - Test 12 (09 Mar 2022)}, paper=1, qno=4, type=multi, subtopic={Properties}, answer={A,B}, answer_src=key, difficulty=7, diff_parts={K1 S4 I1 C1 F1}, diff_note={Reduction formula by parts for I_n then I_2, I_3}]
If $I_n=\displaystyle\int_0^1\dfrac{dx}{(1+x^2)^n}$, $n\in N$, then which of the following statements hold good?
\begin{qoptions}
\opt{A}{$2nI_{n+1}=2^{-n}+(2n-1)I_n$}
\opt{B}{$I_2=\dfrac\pi8+\dfrac14$}
\opt{C}{$I_2=\dfrac\pi8-\dfrac14$}
\opt{D}{$I_3=\dfrac\pi{16}-\dfrac5{48}$}
\end{qoptions}
\end{bankq}
\begin{bankq}[id=SI12-P1-MATH-S1-Q05, source={Score-I - Test 12 (09 Mar 2022)}, paper=1, qno=5, type=multi, subtopic={Leibniz rule}, answer={A,D}, answer_src=key, difficulty=7, diff_parts={K2 S3 I1 C1 F1}, diff_note={Differentiate integral equation to a linear ODE, constant from F(4)=0, then F'(4) and F(8)}]
If $F(x)=\dfrac1{x^2}\displaystyle\int_4^x\left(4t^2-2F'(t)\right)dt$, then $F'(4)$ equals --
\begin{qoptions}
\opt{A}{$\dfrac{32}9$}
\opt{B}{$\dfrac{64}9$}
\opt{C}{$\dfrac{F(8)}{28}$}
\opt{D}{$\dfrac{11F(8)}{28}$}
\end{qoptions}
\end{bankq}
\begin{bankq}[id=SI12-P2-MATH-S2-Q01, source={Score-I - Test 12 (09 Mar 2022)}, paper=2, qno=1, type=single, subtopic={Properties}, answer={3.00}, answer_src=key, difficulty=4, diff_parts={K1 S3 I1 C0 F0}, diff_note={f(n)=n pi/2 via f(n+1)-f(n)=integral of sin((2n+1)t)/sin t}]
If $f(n)=\displaystyle\int_0^{\pi/2}\dfrac{\sin^2n\theta}{\sin^2\theta}\,d\theta$, $(n\in N)$, then the value of $\dfrac{f(15)+f(3)}{f(15)-f(9)}$ is
\end{bankq}
\begin{bankq}[id=SI12-P2-MATH-S2-Q02, source={Score-I - Test 12 (09 Mar 2022)}, paper=2, qno=2, type=single, subtopic={Properties}, answer={2022.00}, answer_src=key, difficulty=5, diff_parts={K2 S3 I1 C0 F0}, diff_note={tan x |x| is odd so vanishes; 2022*2 int0^(1/n) x cos x dx ~ 1/n^2 gives 2022}]
The value of $\displaystyle\lim_{n\to\infty}n^2\int_{-1/n}^{1/n}(2021\tan x+2022\cos x)\,|x|\,dx$
\end{bankq}
\end{document}
